\section{Implementation}\label{implementation}

This chapter reports the concrete technology-dependent choices used to realize the design discussed in \cref{design}. The implementation deliberately favors a lightweight and inspectable stack, so that the distributed behavior of the system remains visible and reproducible without relying on heavy external infrastructure.

\subsection{Protocols and Data Representation}
\label{implementation-protocols}

Two communication protocols are used, each chosen according to the interaction style required by the system.

\textbf{Node-to-node communication.}
Peer communication is implemented over \textbf{TCP}. This choice was preferred over UDP because the project needs ordered and reliable delivery of protocol payloads at the transport level, while still keeping the communication stack simple and under direct control. TCP is therefore used for:
\begin{itemize}
  \item bootstrap and peer discovery traffic;
  \item heartbeat exchange;
  \item membership gossip dissemination;
  \item sensor-state update propagation;
  \item incremental and full synchronization during recovery.
\end{itemize}

Rather than adopting an external RPC or messaging framework, the project uses a custom lightweight application protocol on top of TCP. Payloads are transmitted as length-prefixed frames, which makes message boundaries explicit and keeps the implementation compatible with a persistent stream-oriented transport.

\textbf{Observer communication.}
Human users and external tools interact with the system through a \textbf{read-only HTTP API}. HTTP is appropriate here because the observability plane is request-response oriented, browser-friendly, and naturally suited to dashboards, integration scripts, and test harnesses. Since the observation API does not participate in peer coordination, it is intentionally separated from the internal peer-to-peer protocol.

\textbf{In-transit representation.}
Messages and HTTP responses are represented as \textbf{JSON}. This choice was motivated by readability, ease of debugging, and direct interoperability with browsers and Python tooling. JSON is sufficient for the project because payload sizes are moderate and the main goal is clarity rather than binary compactness. Message validation and dispatch are then handled at the protocol layer after decoding.

The main message families implemented in the peer protocol are:
\begin{itemize}
  \item \texttt{JOIN\_REQUEST} and \texttt{PEER\_LIST} for discovery;
  \item \texttt{PING} and \texttt{PONG} for heartbeat-based liveness observation;
  \item \texttt{GOSSIP\_STATE} for disseminating membership knowledge;
  \item \texttt{SENSOR\_UPDATE} for push-based sensor replication;
  \item \texttt{GET\_DELTA} and \texttt{DELTA\_UNAVAILABLE} for incremental recovery;
  \item \texttt{FULL\_SYNC\_REQUEST} and \texttt{FULL\_SYNC\_RESPONSE} for stronger catch-up after disruption.
\end{itemize}

\subsection{State Management and Storage Choices}
\label{implementation-storage}

The implementation does \textbf{not} use an external database. All relevant information is maintained as in-memory runtime state inside each node:
\begin{itemize}
  \item the current winning sensor records;
  \item replication deltas and sequence metadata;
  \item the peer membership table;
  \item liveness-related information;
  \item topology snapshots;
  \item introspection events and counters.
\end{itemize}

This is consistent with the project scope: the goal is to study replication, convergence, recovery, and observability of current distributed state, not durable long-term business storage. As a consequence, no SQL or NoSQL query layer is required. Queries are replaced by direct access to in-memory state providers, which are then exposed through the HTTP observation endpoints.

From an implementation viewpoint, the most important storage choice is therefore not the selection of a database technology, but the decision to encode merge semantics explicitly in the application logic. Sensor-state convergence is implemented with deterministic Last-Writer-Wins conflict resolution based on timestamp and origin metadata, while membership and liveness views are maintained through dedicated in-memory structures updated by protocol handlers and periodic runtimes.

\subsection{Authentication and Authorization}
\label{implementation-auth}

No explicit authentication or authorization protocol is implemented in the current version of the system. This is a deliberate scope decision, not an omission by accident. The project runs in controlled development and laboratory environments, and the report prioritizes decentralization, consistency, failure handling, and observability over access-control concerns.

Accordingly:
\begin{itemize}
  \item peer-to-peer channels do not currently require node authentication;
  \item the HTTP API does not require login, tokens, or session management;
  \item the system does not implement role-based or attribute-based authorization.
\end{itemize}

The practical effect is that any process able to reach a node in the current deployment can interact with the exposed interface allowed by that channel. This is acceptable for the project scope, but it also marks a clear limitation of the current implementation. Stronger authentication, endpoint protection, and transport hardening are left as future work rather than embedded in the present runtime.

\subsection{Technological details}\label{technological-details}

\textbf{Programming language and runtime.}
The project is implemented in \textbf{Python}. This choice supports rapid experimentation, readable code, and low setup cost for a course project. The implementation makes extensive use of the Python standard library for networking, concurrency, HTTP serving, serialization, logging, and timing, which keeps external dependencies intentionally minimal.

\textbf{Concurrency model.}
The runtime is based on \textbf{threads}, background worker loops, and in-memory queues. This model is used for several long-lived activities that must proceed concurrently:
\begin{itemize}
  \item accepting inbound TCP connections;
  \item maintaining outbound peer workers;
  \item periodic heartbeat and failure-detector evaluation;
  \item periodic push-pull replication;
  \item sensor generation;
  \item HTTP serving for observation.
\end{itemize}

This approach was preferred to a more complex actor framework or asynchronous event-loop architecture because it keeps the execution model explicit and understandable, which is desirable for an educational distributed-systems project.

Concurrent access to local replicated state is handled through a combination of queue-based ingestion and explicit locking. Local sensor providers enqueue normalized sensor events into a thread-safe event queue consumed by the \texttt{NodeStateWorker}; remote protocol handlers can invoke merge operations for \texttt{SENSOR\_UPDATE} and \texttt{FULL\_SYNC\_RESPONSE} messages. The underlying \texttt{NodeStateStore} protects sensor records, update buffers, replication deltas, sequence cursors, and origin-sequence tracking with a mutex. Membership and failure-detector state are protected separately by locks in the peer table and heartbeat monitor. This mitigates in-process race conditions among sensor generation, TCP handlers, replication publishing, heartbeat evaluation, and HTTP polling. The remaining limits are deliberate: state and membership are not updated atomically as one transaction, replication buffers are bounded and volatile, and recovery relies on peer synchronization rather than durable local storage.

\textbf{Configuration management.}
Node configuration is externalized through \textbf{environment variables}. Parameters such as node identity, ports, bootstrap peers, sensor definitions, replication cadence, failure-detector thresholds, and fault-injection options are all configurable from the environment. This makes it easy to instantiate different topologies and experimental scenarios without changing source code.

\textbf{Sensor implementation.}
Sensor sources are implemented as pluggable local producers. The repository includes multiple simulated sensor types, such as numeric, boolean, wave, noise, spike, trend, categorical, incremental, and finite test sensors. This confirms the extensibility objective identified in NFR6 while keeping the whole platform reproducible in the absence of real hardware devices.

\textbf{Networking layer.}
The transport layer is implemented with Python sockets and a custom framed TCP protocol. On the outbound side, each peer connection has its own worker and FIFO send queue; on the inbound side, a threaded TCP server accepts connections, reads framed messages, validates them, and dispatches them to protocol handlers. The implementation also includes bounded worker counts and configurable socket limits, which helps prevent unbounded resource growth in larger experiments.

\textbf{HTTP observation layer.}
The Web API is implemented with Python's `http.server` facilities, in particular a threaded HTTP server and custom request handlers. This is sufficient because the API is read-only and relatively small, so introducing a larger web framework would not provide proportionate benefits for the project scope.

\textbf{Containerization and deployment automation.}
The project uses \textbf{Docker} and \textbf{Docker Compose} to instantiate reproducible node topologies. Multiple compose configurations are provided, including small and medium cluster layouts, so that the same codebase can be exercised under different scales and fault scenarios. Containerization strongly supports NFR5 because it reduces environmental variability across test runs and demonstrations.

\textbf{Testing and quality tools.}
Validation and development are supported by:
\begin{itemize}
  \item \textbf{pytest} for unit and integration testing;
  \item Docker-backed integration harnesses for multi-node and fault-recovery scenarios;
  \item \textbf{pyright} for static type checking;
  \item structured logging and introspection endpoints for runtime diagnosis.
\end{itemize}

\textbf{Dependencies.}
External dependencies are intentionally few. Besides development tools, the main runtime dependency is `python-dotenv`, used to simplify environment-based configuration. This minimal dependency footprint is coherent with the project's emphasis on portability, inspectability, and self-contained deployment.

\subsection{Comparison with Messaging and Streaming Frameworks}
\label{framework-comparison}

One of the learning goals of the project is to compare the implemented peer-to-peer approach with established distributed messaging and streaming technologies. The system intentionally avoids adopting these frameworks as core dependencies because the educational objective is to expose the mechanisms of gossip dissemination, membership maintenance, failure suspicion, and recovery rather than delegating them to middleware.

\textbf{MQTT.}
MQTT is well suited to IoT telemetry because it offers a lightweight publish/subscribe abstraction and a broker-centered communication model. It would simplify sensor ingestion and client distribution, but the broker would become an explicit infrastructural dependency and, unless clustered externally, a possible coordination point. In comparison, Distributed Sensor Hub keeps dissemination decentralized: every node participates directly in replication and membership, which makes failures and convergence behavior more visible for experimentation. The trade-off is that this project must implement discovery, peer communication, and recovery logic itself, while MQTT would provide a simpler application-facing interface.

\textbf{ZeroMQ.}
ZeroMQ provides flexible messaging patterns and efficient socket abstractions without imposing a broker in all configurations. It could reduce the amount of low-level networking code and support several communication topologies. However, it does not by itself define a replicated sensor-state model, membership protocol, Phi Accrual failure detector, or full-state rejoin semantics. Using ZeroMQ would therefore help with transport engineering but would not replace the distributed-system logic that this project wants to study directly. The custom TCP and JSON protocol is less general, but keeps the behavior inspectable and aligned with the course goals.

\textbf{Kafka.}
Kafka provides durable logs, high-throughput ingestion, replay, and strong operational support for stream processing. It is much stronger than this prototype for production-scale persistence, ordering within partitions, and consumer coordination. At the same time, Kafka relies on a managed broker cluster and a log-centric data model. That would shift the project away from autonomous peer replicas toward a centralized streaming substrate. Distributed Sensor Hub instead sacrifices durable history and large-scale throughput in favor of node autonomy, local availability during partitions, and direct observation of eventual consistency through timestamp-based merging.

\textbf{Redis Streams.}
Redis Streams offers a compact and practical stream abstraction with consumer groups and persistence options. It would be useful for buffering sensor events and coordinating consumers in a small deployment, but it still introduces a server-side data store as a central component unless Redis is deployed and managed in a replicated configuration. The implemented system does not need an external store for normal operation: each node keeps its own runtime replica and repairs it through gossip, deltas, and full synchronization. This improves conceptual decentralization, while Redis Streams would provide stronger operational convenience and better persistence primitives.

Overall, the implemented architecture is not meant to outperform these frameworks in production maturity. Its value is educational and experimental: it makes consistency, scalability, and fault-tolerance trade-offs explicit. MQTT and Redis Streams would simplify integration but introduce broker or server dependencies; ZeroMQ would improve transport flexibility but leave higher-level replication semantics to the application; Kafka would provide durable scalable streaming but change the architecture into a brokered log-based system. Distributed Sensor Hub instead prioritizes decentralized membership, anti-entropy replication, adaptive failure detection, and recovery after partitions, accepting simpler storage and course-project scale as deliberate constraints.

\clearpage
